\chapter{Solução Proposta} \label{Cap:Metodologia}

\section{Considerações iniciais}
\label{sec:sol_consideracoes_iniciais}

Este capítulo detalha a concepção arquitetural, os modelos matemáticos e a implementação prática da solução desenvolvida para a extração e a classificação de transações de faturas de cartão de crédito: formaliza as questões de pesquisa e as hipóteses norteadoras, descreve a metodologia baseada no CRISP-DM adaptado, caracteriza a fonte de dados e a amostragem, especifica os microsserviços do ecossistema \textit{data lakehouse}, pormenoriza a engenharia do agente inteligente em grafo (LangGraph), estabelece as salvaguardas de conformidade com a LGPD e delineia a estratégia de garantia de qualidade e testes de software.

\section{Questões de pesquisa e hipótese central}
\label{sec:sol_questoes_pesquisa}

A formulação e a avaliação experimental deste trabalho são guiadas por cinco Questões de Pesquisa (Q1 a Q5) fundamentadas na literatura:

\begin{itemize}
	\item \textbf{Q1:} Quais técnicas de representação textual (esparsas versus densas) e algoritmos de aprendizado (estatísticos clássicos versus modelos neurais baseados em subpalavras e agentes LLM) apresentam maior eficácia para categorizar transações a partir de descrições ruidosas e despadronizadas?
	\item \textbf{Q2:} Qual é o nível de acurácia global e F1-Macro alcançável por um classificador automático de gastos operando sobre faturas reais e ruidosas, considerando uma taxonomia de 10 a 20 categorias de despesas?
	\item \textbf{Q3:} Quais são as principais naturezas de erro cometidas pelos classificadores automatizados (ambiguidades de \textit{marketplace}, erros de leitura de OCR, nomes informais de lojistas) e de que forma mecanismos como \textit{Chain-of-Thought} e busca na web podem mitigá-los?
	\item \textbf{Q4:} Em que medida a injeção de metadados contextuais (origem do arquivo digital versus óptico, histórico de compras do usuário e descarte de prefixos de \textit{gateways}) contribui para o incremento do desempenho preditivo?
	\item \textbf{Q5:} Qual é a viabilidade operacional e o custo (latência, rastreabilidade e acurácia condicional do ramo de raciocínio) de um ciclo de vida de MLOps com aprendizado ativo e intervenção humana (HITL) para manter e atualizar o sistema, executado integralmente em infraestrutura local? A caracterização empírica de \textit{data drift} e \textit{concept drift} ao longo do tempo (Seção~\ref{sec:fund_drift}), embora motivadora do ciclo de MLOps, não é objeto de medição nesta pesquisa (Seção~\ref{sec:concl_limitacoes}).
\end{itemize}

A \textbf{hipótese central} submetida a teste nesta pesquisa estabelece que:
\begin{quote}
	\textit{``Uma esteira híbrida que integra normalização léxica de gateways, busca por similaridade semântica em banco vetorial (pgvector), raciocínio dinâmico via agente LLM com acesso a ferramentas de busca e intervenção humana programada (\textit{Human-in-the-Loop}) é capaz de classificar transações de cartão de crédito com acurácia global e F1-Macro superiores a 85\% em uma taxonomia de 14 categorias.''}
\end{quote}
A conformidade com a LGPD e a execução em infraestrutura local, embora indissociáveis da solução, não integram a hipótese: são requisitos não funcionais do projeto (Capítulo~\ref{Cap:Introdução}), verificados por inspeção arquitetural (Seção~\ref{sec:sol_lgpd}) e não submetidos a teste de hipótese, por não constituírem proposições falsificáveis por experimento. Por conjugar dois critérios quantitativos distintos sob uma mesma afirmação, a hipótese central é decomposta em duas sub-hipóteses independentes, cada uma submetida a verificação e veredito próprios (Seção~\ref{sec:val_respostas_questoes}):
\begin{itemize}
	\item \textbf{H1a (Acurácia):} a acurácia global do classificador é superior a 85\%;
	\item \textbf{H1b (F1-Macro):} o F1-Macro do classificador é superior a 85\%.
\end{itemize}
Em termos de teste estatístico, cada sub-hipótese é a hipótese alternativa de um par formalizado a priori (procedimentos na Seção~\ref{sec:fund_estatistica}): para a acurácia, $H_{0a}$: acurácia populacional $\le 85\%$ contra $H_{1a}$: acurácia populacional $> 85\%$, verificado por teste binomial exato unilateral (Equação~\ref{eq:binomial}) e pelo intervalo de confiança de Wilson (Equação~\ref{eq:wilson}) sobre a proporção de acertos; para o F1-Macro, $H_{0b}$: F1-Macro populacional $\le 85\%$ contra $H_{1b}$: F1-Macro populacional $> 85\%$, para o qual não há teste exato e adota-se a regra operacional de rejeitar $H_{0b}$ somente se o limite inferior do intervalo de confiança de 95\% obtido por reamostragem \textit{bootstrap} for superior a 85\%. Adota-se o nível de significância $\alpha = 0{,}05$; uma sub-hipótese é considerada sustentada somente se a estimativa pontual supera o limiar \textit{e} a hipótese nula correspondente é rejeitada pelo critério declarado. Os procedimentos são detalhados na Seção~\ref{sec:val_significancia}. Como a formulação original é conjuntiva (\textit{E lógico} entre H1a e H1b), a hipótese central como um todo é considerada refutada caso qualquer uma das sub-hipóteses seja refutada, independentemente do resultado da outra. O \textbf{objeto} da hipótese é a esteira híbrida enunciada -- isto é, o agente em grafo de estados descrito na Seção~\ref{sec:sol_agente_inteligente} --, e é sobre ela que as duas sub-hipóteses são verificadas (Seção~\ref{sec:val_respostas_questoes}). As demais arquiteturas avaliadas (linhas de base esparsas, modelo no nível de byte e busca vetorial isolada) cumprem o papel de referência comparativa exigido pela Q1; seus resultados são reportados frente ao mesmo limiar de 85\%, mas não substituem a esteira híbrida como objeto da verificação.

\section{Metodologia da pesquisa e arcabouço CRISP-DM adaptado}
\label{sec:sol_crisp_dm}

A pesquisa classifica-se, quanto à sua natureza, como \textbf{aplicada}, voltada à resolução do problema prático de gestão financeira pessoal \cite{gerhardt2009metodos}. Quanto à abordagem e procedimentos, define-se como uma pesquisa \textbf{experimental e quantitativa} \cite{gil2019metodos}, estruturada no controle de variáveis em ambiente controlado e na mensuração estatística de métricas preditivas.

O ciclo metodológico adotou o arcabouço padrão da indústria \textit{Cross-Industry Standard Process for Data Mining} (CRISP-DM) \cite{chapman2000crisp, wirth2000crisp}, estendido com diretrizes contemporâneas de MLOps \cite{kreuzberger2022mlops} conforme as seis etapas sequenciais:
\begin{enumerate}
	\item \textbf{Compreensão do negócio:} Mapeamento das regras de faturas de cartão de crédito no Brasil, identificação de intermediadores de pagamento e taxonomia de categorias orçamentárias;
	\item \textbf{Compreensão dos dados:} Coleta e exploração de conjuntos de dados públicos oficiais (CPGF) e de faturas reais, com as transações pseudonimizadas na base (Seção~\ref{sec:sol_lgpd}), com inspeção qualitativa do ruído nativo introduzido por intermediadores de pagamento e terminais de captura;
	\item \textbf{Preparação dos dados:} Engenharia de extração com preservação de leiaute (Docling V2) e comparação experimental com reconhecimento óptico (Tesseract OCR), normalização léxica com expressões regulares, descarte de jargões corporativos e cálculo de \textit{embeddings} vetoriais;
	\item \textbf{Modelagem:} Desenvolvimento comparativo englobando linhas de base clássicas (TF-IDF + Random Forest / Naive Bayes), ajuste fino do ByT5 e construção do agente inteligente em grafo de estados (LangGraph com pgvector e LLM);
	\item \textbf{Avaliação:} Mensuração de métricas estatísticas (Acurácia, Precisão, Revocação, F1-Macro e Matriz de Confusão), estudos de ablação e auditoria de latência;
	\item \textbf{Implantação e ciclo de vida:} Conteinerização em microsserviços via Docker, armazenamento em camadas com MinIO e mensageria com Kafka, rastreamento de experimentos com MLflow \cite{zaharia2018mlflow} e autenticação da aplicação \textit{web} via Logto.
\end{enumerate}

\section{Fonte de dados, amostragem e natureza dos dados}
\label{sec:sol_fonte_dados}

A pesquisa opera exclusivamente sobre \textbf{dados reais}, e não sintéticos, provenientes de dois acervos de naturezas distintas e papéis complementares, caracterizados em detalhe na Seção~\ref{sec:val_datasets}:
\begin{itemize}
	\item \textbf{Acervo governamental CPGF} (Cartão de Pagamento do Governo Federal): 477.740 transações públicas obtidas no Portal da Transparência \cite{brasilcpgf}, empregadas apenas para a caracterização exploratória do ruído textual em escala, e não para a mensuração preditiva;
	\item \textbf{Acervo pessoal de faturas de cartão de crédito} de dois emissores brasileiros (C6 Bank e Nubank), pertencentes a dois portadores de cartão: 49 faturas do C6 Bank em PDF (357 páginas) pareadas com os 49 arquivos CSV oficiais correspondentes, que constituem a verdade de referência para a medição da qualidade da extração; e os lançamentos em CSV de ambos os emissores, que alimentam a base rotulada de classificação.
\end{itemize}

A decisão de avaliar a classificação diretamente sobre o \textbf{ruído real} das faturas -- truncamentos, prefixos de intermediadores de pagamento, sufixos de autorização e nomes informais aglutinados -- em vez de gerá-lo sinteticamente é uma escolha metodológica deliberada: ruído sintético reproduz apenas os fenômenos que o pesquisador já conhece, ao passo que o acervo real contém a distribuição efetiva desses fenômenos e aqueles ainda não catalogados. O custo dessa escolha é o porte reduzido do acervo e sua origem em apenas dois portadores, limitação de validade externa registrada na Seção~\ref{sec:concl_limitacoes}.

Quanto à \textbf{amostragem}, a base rotulada de classificação foi construída a partir dos lançamentos dos dois emissores após deduplicação e limitação a no máximo 12 lançamentos por estabelecimento -- de modo que a avaliação meça a amplitude de comércios distintos, e não a frequência de poucos estabelecimentos dominantes --, totalizando 1.660 transações em 11 categorias efetivamente povoadas. A \textbf{rotulagem} é semiautomática, em duas camadas: um dicionário léxico de alta precisão validado manualmente sobre amostra (74\% dos rótulos) e, nos casos não cobertos, o mapeamento da categoria nativa do emissor para a taxonomia canônica. O particionamento em treino, validação e teste cego é feito por \textbf{descrição normalizada}: nenhuma descrição do teste aparece no treino, embora um mesmo comércio possa figurar nos dois conjuntos quando sua grafia varia entre lançamentos --- limite quantificado na Seção~\ref{sec:val_dataset_classificacao}. O detalhamento completo -- composição por emissor, distribuição por categoria, cobertura da taxonomia e risco de circularidade entre rotulagem e desempenho -- é apresentado na Seção~\ref{sec:val_dataset_classificacao}.

\section{Arquitetura geral da solução e ecossistema de dados em camadas}
\label{sec:sol_arquitetura}

A arquitetura geral do sistema foi concebida sob o paradigma orientado a eventos e microsserviços desacoplados, implementada com tecnologias de código aberto (\textit{Open Source}) orquestradas via Docker Compose. Trata-se de uma arquitetura em camadas \textbf{inspirada no paradigma \textit{lakehouse}} \cite{armbrust2021lakehouse}. Desse paradigma, esta pesquisa implementa a camada Bronze em armazenamento de objetos (arquivos brutos imutáveis) e a camada de dados tratados em banco relacional e vetorial. As propriedades que caracterizam um \textit{lakehouse} completo -- formato de tabela aberto com transações ACID sobre o armazenamento de objetos e consumo analítico desacoplado -- são apontadas como evolução no Capítulo~\ref{Cap:Conclusoes}. O termo é mantido nas figuras e nos nomes de componentes por concisão, com essa ressalva. A Figura~\ref{fig:arquitetura_mlops} ilustra a topologia da infraestrutura efetivamente implementada e exercitada neste trabalho.

\begin{figure}[htb]
	\centering
	\caption{Arquitetura em camadas inspirada no paradigma \textit{lakehouse}, orientada a eventos, com MLOps e agentes de IA.}
	\label{fig:arquitetura_mlops}
	\vspace{0.3cm}
	
	\resizebox{0.95\textwidth}{!}{%
		\begin{tikzpicture}[
			box/.style={draw, rectangle, rounded corners, align=center, minimum width=3.2cm, minimum height=1.4cm, font=\sffamily\small, fill=white},
			storage/.style={draw, cylinder, shape border rotate=90, aspect=0.25, align=center, minimum width=2.8cm, minimum height=1.6cm, font=\sffamily\scriptsize, fill=white},
			arrow/.style={-Latex, thick, draw=gray!80!black},
			dashed_arrow/.style={Latex-, dashed, thick, draw=gray!80!black},
			elabel/.style={font=\sffamily\scriptsize, fill=white, inner sep=1.5pt}
			]

			% Linha 1: Ingestão e Usuários
			\node[box, fill=gray!10]   (user)  at (-5.6, 0)  {Usuário / Cliente \\ (Envio da Fatura)};
			\node[box, fill=blue!10]   (api)   at (0, 0)     {API WebFlux \\ (Ingestão Reativa)};
			\node[box, fill=orange!10] (kafka) at (5.6, 0)   {Apache Kafka \\ (Mensageria e Tópicos)};
			\node[box, fill=purple!15] (ds)    at (11.2, 0)  {Cientista de Dados \\ (Jupyter / MLflow)};

			% Linha 2: Workers de Extração e IA
			\node[storage, fill=red!10]   (bronze_raw) at (-5.6, -3.6) {MinIO Bronze \\ (PDFs / CSVs Brutos)};
			\node[box, fill=green!10]     (ocr)        at (0, -3.6)     {Trabalhador de Extração \\ (Docling V2)};
			\node[box, fill=purple!10]    (byt5)       at (5.6, -3.6)   {Agente Inteligente \\ (LangGraph + LLM)};
			\node[storage, fill=orange!20](mlops)      at (11.2, -3.6)  {MinIO MLOps \\ (Artefatos MLflow)};

			% Linha 3: Camada de dados tratados
			\node[storage, fill=teal!10] (silver)     at (5.6, -7.6) {PostgreSQL + pgvector \\ (Faturas, Transações, \\ Base de Conhecimento e \textit{Embeddings})};

			% (Camada de consumo analítico -- prevista para evolução, ver Cap. 6)

			% Fundo Agrupador Data Lakehouse
			\begin{scope}[on background layer]
				\node[fill=gray!5, rounded corners, draw=gray, dashed, fit=(bronze_raw) (silver) (mlops), inner sep=18pt] (datalake) {};
				\node[anchor=north, font=\sffamily\bfseries\scriptsize, align=center, text=gray!70!black, yshift=-2pt] at (datalake.south) {Camadas de dados (armazenamento de objetos S3 e banco relacional/vetorial)};
			\end{scope}

			% Conexões / Setas
			\draw[arrow] (user) -- (api);
			\draw[arrow] (api) -- node[elabel] {Publica Evento} (kafka);
			\draw[arrow] (api.south) -- node[elabel, pos=0.3] {Persiste Bruto} (bronze_raw.north);

			\draw[arrow] (kafka.south west) -- node[elabel, pos=0.26] {Consome Evento} (ocr.north east);
			\draw[dashed_arrow] (ocr.west) -- node[elabel] {Lê Arquivo} (bronze_raw.east);
			\draw[arrow] (ocr.south) |- node[elabel, align=center, pos=0.22, right] {Grava Fatura e\\Transações (pendentes)} (silver.west);

			\draw[arrow] (ds.south west) -- node[elabel, sloped, pos=0.3] {Executa em lote} (byt5.north east);
			\draw[dashed_arrow] ([xshift=-0.9cm]byt5.south) -- node[elabel, left, pos=0.5] {Lê Pendentes / Busca Vetorial} ([xshift=-0.9cm]silver.north);
			\draw[arrow] ([xshift=0.9cm]byt5.south) -- node[elabel, right, pos=0.5] {Grava Classificação} ([xshift=0.9cm]silver.north);

			\draw[arrow] (ds.south) -- node[elabel, pos=0.32] {Rastreia Runs} (mlops.north);
			\draw[dashed_arrow] (byt5.east) -- node[elabel] {Registra métricas} (mlops.west);

		\end{tikzpicture}
	}

	\vspace{0.2cm}
	{\small \textbf{Fonte:} Elaborada pelo autor (\the\year). \textbf{Nota:} o agente não é acionado por evento do Kafka; ele é executado em lote (\codepath{notebooks/experimentos/05_agente.py}) sobre as transações persistidas com situação pendente pelo trabalhador de extração, que grava diretamente no PostgreSQL. O \textit{bucket} intermediário de textos extraídos previsto no provisionamento (\texttt{bronze-txt}) não é utilizado pela esteira.}
\end{figure}

O ecossistema divide-se nos seguintes serviços especializados:
\begin{itemize}
	\item \textbf{MinIO (armazenamento de objetos S3):} Provedor de armazenamento em nuvem local compatível com a API AWS S3. A esteira utiliza o \textit{bucket} \texttt{bronze-raw}, particionado por fonte e data (\codepath{bronze/source=}\texttt{\{c6|nubank\}}\allowbreak\codepath{/year=/month=/day=/}\allowbreak\texttt{\{uuid\}}\allowbreak\codepath{/original.}\texttt{\{pdf|csv\}}), para a preservação imutável dos arquivos brutos, e o \textit{bucket} \texttt{mlflow-artifacts} como raiz de artefatos do servidor MLflow. O \textit{script} de provisionamento cria ainda os \textit{buckets} \texttt{bronze-txt} e \texttt{silver}, reservados para uma camada de textos extraídos e para tabelas abertas que não foram materializadas nesta pesquisa: os dados extraídos são gravados diretamente no PostgreSQL pelo trabalhador de extração;
	\item \textbf{Apache Kafka:} Barramento de mensageria que desacopla temporalmente a camada de ingestão do trabalhador de extração por meio do tópico \texttt{datalake-\allowbreak bronze-\allowbreak ingestao-\allowbreak topic}. O consumidor opera com uma partição, confirmação automática de deslocamento (\textit{auto-commit}) e leitura a partir das mensagens mais recentes, o que configura semântica de entrega \textit{no máximo uma vez}: um evento consumido durante uma falha do trabalhador não é reprocessado -- comportamento aceitável para o protótipo e registrado como evolução no Capítulo~\ref{Cap:Conclusoes};
	\item \textbf{PostgreSQL 16 com extensão pgvector:} Banco de dados relacional e vetorial encarregado de armazenar usuários, faturas, transações, a base de conhecimento de estabelecimentos e os vetores de características densas de 384 dimensões gerados pelo modelo \modelemb{}. A base vetorial efetivamente consultada pelo agente é a coleção \texttt{baseconhecimento}, mantida pela integração \texttt{langchain-postgres} nas tabelas \texttt{langchain\_pg\_collection} e \texttt{langchain\_pg\_embedding}; a tabela relacional \texttt{estabelecimentos} (Apêndice~\ref{apendice:esquema_banco}) guarda cópia do mesmo \textit{embedding} para consulta relacional e é povoada na mesma etapa (Seção~\ref{sec:val_protocolo});
	\item \textbf{Servidor MLflow:} Servidor centralizado de MLOps apoiado por um banco relacional PostgreSQL próprio (\texttt{mlops\_db}) e por um \textit{bucket} S3. A esteira experimental registra, por chamadas explícitas ao cliente MLflow, os parâmetros, as métricas globais, os relatórios por categoria, as predições e as figuras de cada arquitetura sob o experimento \texttt{experimento\_faturas}. O registro automático (\textit{autologging}) das chamadas LangChain/LangGraph -- com \textit{prompt}, \textit{tokens} e cadeia de pensamento por transação -- foi empregado apenas no \textit{notebook} exploratório do agente e não integra a esteira reprodutível que gera os resultados do Capítulo~\ref{Cap:Ava_Exp};
	\item \textbf{Logto IAM:} Servidor de identidade baseado em OpenID Connect (OIDC) e OAuth 2.0 que autentica o usuário da aplicação \textit{web}. Nesta versão, os pontos de acesso da API de ingestão e da API de classificação não validam o \textit{token} emitido pelo Logto, e todos os registros são associados a um único usuário interno; o controle de acesso por papéis e o isolamento multi-inquilino são, portanto, previstos pela arquitetura mas não implementados (Seções~\ref{sec:sol_lgpd} e~\ref{sec:concl_limitacoes}).
\end{itemize}

A Tabela~\ref{tab:tecnologias} consolida, para cada componente da solução, a linguagem, o arcabouço principal e a versão efetivamente utilizada; o ambiente de execução dos experimentos é descrito na Seção~\ref{sec:val_ambiente}.

\begin{table}[htb]
	\centering
	\caption{Componentes da solução: linguagem, arcabouço principal e versão efetivamente utilizada.}
	\label{tab:tecnologias}
	\small
	\begin{tabularx}{\textwidth}{>{\raggedright\arraybackslash}p{4.4cm} p{2.1cm} X}
		\toprule
		\textbf{Componente} & \textbf{Linguagem} & \textbf{Arcabouço principal e versão} \\
		\midrule
		API de ingestão reativa (\codepath{service-ingestao}) & Java 25 & Spring Boot 4.0.6 (WebFlux, Spring Kafka, Reactor Kafka 1.3.22); AWS SDK for Java v2 2.27.21 (cliente S3 assíncrono contra o MinIO); ICU4J 78.3 (detecção de codificação) \\
		Trabalhador de extração (\codepath{worker-prata}) & Python 3.11 & IBM Docling 2.112.0 sobre PyPDFium2 5.9.0, com OCR desativado (entradas nativas digitais); kafka-python; SQLAlchemy \\
		Agente de classificação (\codepath{notebooks/modules}) & Python 3.11 & LangGraph 1.2.8; LangChain 1.3.12; Pydantic 2.13.4; langchain-postgres 0.0.17; RapidFuzz 3.14.5 \\
		API de classificação (\codepath{service-classification}) & Python 3.11 & FastAPI 0.138.0; protótipo de classificação \textit{zero-shot} por similaridade com \codepath{intfloat/multilingual-e5-large}, não integrado ao agente nem avaliado no Capítulo~\ref{Cap:Ava_Exp} \\
		Esteira experimental (\codepath{notebooks/experimentos} e \textit{notebook} de avaliação de OCR) & Python 3.11 & scikit-learn 1.7.2; PyTorch 2.12.1 (CUDA 13.0); Transformers 5.12.1; Sentence-Transformers 5.6.0; PyTesseract 0.3.13 sobre Tesseract 5.5.0; JiWER 4.0.0; cliente MLflow 3.14.0 \\
		Modelo de \textit{embeddings} & -- & \modelemb{} (384 dimensões) \\
		Modelo de raciocínio local & -- & \codepath{google/gemma-4-e4b} servido pelo LM Studio (API compatível com OpenAI) \\
		Banco relacional e vetorial & SQL & PostgreSQL 16 com pgvector (imagem \codepath{pgvector/pgvector:pg16}) \\
		Barramento de eventos & -- & Apache Kafka 4.3.0 \\
		Armazenamento de objetos & -- & MinIO (imagem \texttt{minio/minio}, API S3) \\
		Servidor de rastreamento & -- & MLflow (imagem \codepath{ghcr.io/mlflow/mlflow}) \\
		Identidade e acesso & -- & Logto (imagem \texttt{svhd/logto}), OIDC/OAuth 2.0 \\
		Aplicação web (\texttt{web}, submódulo) & TypeScript & Next.js 16.2.10; React 19.2.4; \texttt{@logto/next} 4.2.10 \\
		Orquestração de contêineres & -- & Docker Engine 29.4.0 com Docker Compose \\
		\bottomrule
	\end{tabularx}
	\fonte{Elaborada pelo autor (\the\year), a partir de \codepath{service-ingestao/build.gradle}, \codepath{web/package.json}, \codepath{infra/docker-compose-externalinfra.yml} e do ambiente Python que executou os experimentos (versões instaladas conferidas em setembro de 2026). Versões das imagens marcadas como \texttt{latest} no \textit{compose} não são fixadas pelo repositório. O ambiente de execução (máquinas e sistema operacional) é descrito na Seção~\ref{sec:val_ambiente}.}
\end{table}

A camada de consumo analítico dedicada (motor SQL distribuído e painéis gerenciais de \textit{Business Intelligence}) e a orquestração de esteiras em lote não integram o escopo implementado e avaliado nesta monografia, sendo apontadas como evolução natural da solução no Capítulo~\ref{Cap:Conclusoes}. A aplicação \textit{web} (submódulo \texttt{web}) encontra-se em estágio de protótipo: implementa a autenticação via Logto e a interface de envio de faturas, mas ainda não está conectada à API de ingestão nem ao PostgreSQL; a visualização orçamentária e a tela de confirmação humana são, por isso, entregas previstas e não realizadas nesta versão (Seção~\ref{sec:concl_trabalhos_futuros}).

\section{Esteira de ingestão e extração híbrida de documentos}
\label{sec:sol_extracao_hibrida}

A camada de extração foi concebida para processar documentos digitais nativos (CSV e PDF com camada de texto) e, arquiteturalmente, arquivos ópticos escaneados, por meio de um fluxo em três etapas:

\begin{enumerate}
	\item \textbf{Identificação de metadados no \textit{data lakehouse}:} A API de ingestão persiste cada arquivo em \texttt{bronze-raw} sob um caminho particionado (ex.: \codepath{bronze/source=nubank/year=2026/month=08/day=09/}\allowbreak\texttt{\{uuid\}}\codepath{/original.csv}) e o anota com metadados S3 (emissor de origem, formato detectado, \textit{hash} SHA-256, codificação de caracteres e \textit{flag} \texttt{requer-ocr}). Somente os formatos CSV e PDF são aceitos pela ingestão; OFX e JSON, citados no escopo do Capítulo~\ref{Cap:Introdução}, constam apenas da heurística de metadados e não têm analisador implementado. No agente, a função determinística \codepath{extrair_metadados_datalake} inspeciona esse caminho e deriva o banco de origem, as partições temporais e a \textit{flag} \texttt{is\_digital\_nativo}, que alterna a instrução de origem do dado no \textit{prompt} do modelo de linguagem (Apêndice~\ref{apendice:system_prompt}). A heurística classifica como nativos digitais os arquivos \texttt{.csv}, \texttt{.ofx}, \texttt{.json}, \texttt{.txt} e \texttt{.xml}; qualquer outra extensão -- inclusive \texttt{.pdf}, que no acervo desta pesquisa é sempre nativo digital -- recebe a instrução de considerar possíveis erros de OCR. Essa imprecisão, e o fato de o metadado ter permanecido fixo durante o teste cego, são discutidos na Seção~\ref{sec:val_ambiente};
	\item \textbf{Conversão com Docling V2 (leiaute tabular):} Em faturas em formato PDF (Nubank e C6 Bank), o trabalhador de extração (\codepath{worker-prata}) utiliza a classe \texttt{Document\allowbreak Converter} da biblioteca IBM Docling V2 com o mecanismo de leitura PyPdfium para a camada de texto nativa e com os modelos neurais de análise de leiaute e reconstrução de tabelas (TableFormer, \citeonline{nassar2022tableformer}) executados em GPU (CUDA). O reconhecimento óptico interno do Docling é mantido \textbf{desativado} (\texttt{do\_ocr = False}), pois as entradas são nativas digitais e protegidas por senha. O conversor mapeia itens de tabela com caixas delimitadoras e exporta o conteúdo estruturado em Markdown delimitado por barras verticais (\texttt{|}), preservando a correspondência exata entre data, final do cartão de crédito, descrição do estabelecimento e valor em reais (R\$). Analisadores específicos por emissor (\codepath{extrator_c6bank}, \codepath{extrator_nubank}) convertem as linhas em transações e as persistem no PostgreSQL com situação pendente de classificação;
	\item \textbf{Extração com Tesseract OCR e filtro de confiança:} Para recibos em imagem estrita ou páginas sem camada de texto, foi implementada a rotina \codepath{extrair_texto_avancado} (\codepath{notebooks/modules/utils.py}), que rasteriza as páginas a 300 DPI e aciona o Tesseract OCR em modo OEM 3 (redes LSTM) e PSM 6 (bloco uniforme de texto), aplicando à saída de \texttt{image\_to\_data} um \textbf{filtro de confiança} $>$~40\% que descarta fragmentos visuais irrecuperáveis.
\end{enumerate}

Cabe explicitar o estado de integração da terceira etapa: a rotina Tesseract está implementada como componente autônomo e foi exercitada na bateria comparativa de extração (Seção~\ref{sec:val_resultados_ocr}), mas \textbf{não está acoplada ao trabalhador de extração} -- não há, nesta versão, roteamento automático que a acione quando o Docling não identifica tabelas ou quando o arquivo é uma imagem. Tampouco foi submetida a validação sobre documentos genuinamente escaneados ou fotografados: a avaliação quantitativa do Capítulo~\ref{Cap:Ava_Exp} opera inteiramente sobre PDFs nativos digitais do C6 Bank, artificialmente rasterizados para viabilizar a comparação com o Docling V2. Ambas as lacunas são registradas como limitação do trabalho (Seção~\ref{sec:concl_limitacoes}).

\section{Engenharia do agente de classificação inteligente}
\label{sec:sol_agente_inteligente}

O núcleo de tomada de decisão do sistema foi codificado como um grafo de estados direcionado (StateGraph) no arcabouço LangGraph, suportado pelo estado tipado \texttt{AgentState} e ilustrado na Figura~\ref{fig:stategraph}. O fluxo operacional do agente é estruturado nas seguintes etapas coordenadas:

\begin{figure}[htb]
	\centering
	\caption{Grafo de estados do agente de classificação (StateGraph no LangGraph), com a condição de roteamento de cada aresta, conforme codificado na função \texttt{build\_graph}.}
	\label{fig:stategraph}
	\vspace{0.3cm}
	\resizebox{0.92\textwidth}{!}{%
	\begin{tikzpicture}[
		n/.style={draw, rounded corners, align=center, font=\ttfamily\scriptsize, minimum width=3.6cm, minimum height=0.8cm, fill=white},
		term/.style={draw, rounded corners, align=center, font=\sffamily\scriptsize, minimum width=1.6cm, minimum height=0.6cm, fill=gray!15},
		vec/.style={n, fill=teal!10}, llm/.style={n, fill=purple!10}, hitl/.style={n, fill=orange!20}, db/.style={n, fill=blue!8},
		a/.style={-Latex, thick, draw=gray!70!black},
		lab/.style={font=\sffamily\scriptsize, fill=white, inner sep=1.5pt, align=center}
		]
		\node[term] (start) at (0, 0) {início};
		\node[n]    (norm)  at (0, -1.5) {normalizar};
		\node[n]    (hist)  at (-5.8, -3.1) {buscar\_historico\_marketplace};
		\node[vec]  (sim)   at (0, -4.8) {buscar\_por\_similaridade};
		\node[llm]  (llm)   at (4.6, -7.0) {classificar\_com\_llm};
		\node[llm]  (tools) at (11.8, -7.0) {tools\\{\sffamily(busca na web)}};
		\node[llm]  (estr)  at (4.6, -9.2) {estruturar\_saida\_llm};
		\node[hitl] (hitl)  at (11.8, -11.2) {aguardar\_confirmacao\\{\sffamily(HITL, interrupt)}};
		\node[db]   (vs)    at (4.6, -13.2) {salvar\_vectorstore};
		\node[db]   (res)   at (0, -15.4) {salvar\_resultado};
		\node[term] (end)   at (0, -16.8) {fim};

		\draw[a] (start) -- (norm);
		\draw[a] (norm.west) -| node[lab, pos=0.45, above] {eh\_marketplace verdadeiro} (hist.north);
		\draw[a] (hist.south) |- (sim.west);
		\draw[a] (norm) -- node[lab, right] {caso contrário} (sim);
		\draw[a] ([xshift=-1.0cm]sim.south) -- node[lab, left, pos=0.5] {similaridade $\ge$ 0,85} ([xshift=-1.0cm]res.north);
		\draw[a] (sim.east) -| node[lab, pos=0.25, above] {similaridade $<$ 0,85} (llm.north);
		\draw[a] ([yshift=4pt]llm.east) -- node[lab, above] {chamada de ferramenta} ([yshift=4pt]tools.west);
		\draw[a] ([yshift=-4pt]tools.west) -- node[lab, below] {evidências recuperadas} ([yshift=-4pt]llm.east);
		\draw[a] (llm) -- node[lab, right] {caso contrário} (estr);
		\draw[a] (estr.east) -| node[lab, pos=0.35, above] {requer\_confirmacao\\ou confiança $\le$ 0,80} (hitl.north);
		\draw[a] (hitl.south) |- node[lab, pos=0.7, above] {retomada (Command resume)} (vs.east);
		\draw[a] (estr) -- node[lab, right] {caso contrário} (vs);
		\draw[a] (vs.west) -| node[lab, pos=0.8, right] {persiste somente se\\confiança $>$ 0,85} ([xshift=1.0cm]res.north);
		\draw[a] (res) -- (end);
	\end{tikzpicture}
	}
	\fonte{Elaborada pelo autor (\the\year), a partir de \codepath{notebooks/experimentos/05_agente.py} (função \texttt{build\_graph}).}
	\begin{flushleft}
		\footnotesize \textbf{Nota:} correspondência entre identificadores e função -- \texttt{normalizar}: limpeza léxica; \texttt{buscar\_historico\_marketplace}: consulta ao histórico do usuário; \texttt{buscar\_por\_similaridade}: busca vetorial por similaridade (pgvector); \texttt{classificar\_com\_llm}: inferência com modelo de linguagem; \texttt{tools}: pesquisa autônoma na web; \texttt{estruturar\_saida\_llm}: validação estruturada da resposta do LLM; \texttt{aguardar\_confirmacao}: confirmação humana (HITL); \texttt{salvar\_vectorstore}/\texttt{salvar\_resultado}: persistência no banco de dados. A confirmação humana é alcançada somente após a estruturação da saída do LLM. Descrição de cada nó na lista da Seção~\ref{sec:sol_agente_inteligente}; arestas no Quadro~\ref{quadro:arestas_agente}.
	\end{flushleft}
\end{figure}

\begin{enumerate}
	\item \textbf{Nó \texttt{normalizar}:} Executa a limpeza léxica da descrição da transação com Regex, eliminando duplicações em torno de asteriscos (ex.: ``TIM*TIM'' $\to$ ``TIM''), descartando termos de pagamento (``PAG*'', ``PIX'', ``COMPRA'') e sufixos societários (``LTDA'', ``ME''). Realiza a resolução de aliases de marcas e verifica se a compra pertence à lista de \textit{marketplaces} ambíguos (\texttt{eh\_marketplace}) -- as expressões regulares completas, a lista de \textit{marketplaces} e o dicionário de aliases encontram-se reproduzidos no Apêndice~\ref{apendice:normalizacao}, acompanhados de uma análise crítica do artefato;
	\item \textbf{Roteamento pós-normalização (\texttt{rota\_apos\_normalizar}):} Se a transação for detectada como \textit{marketplace}, o fluxo passa primeiro pelo nó \texttt{buscar\_historico\_marketplace} antes de prosseguir; caso contrário, é direcionado diretamente ao nó de busca vetorial \texttt{buscar\_por\_similaridade}. Na versão avaliada, a detecção compara o nome canônico devolvido pela resolução de aliases com as chaves do dicionário de \textit{marketplaces}; como o alias canônico do Mercado Livre (``mercado livre'') difere da chave correspondente (``mercado\_livre''), esse \textit{marketplace} -- o mais frequente do acervo -- nunca foi roteado pelo ramo de histórico (Apêndice~\ref{apendice:normalizacao}, Seção~\ref{sec:apendice_normalizacao_critica});
	\item \textbf{Nó \texttt{buscar\_historico\_marketplace}:} Consulta o histórico consolidado do usuário no banco PostgreSQL (\texttt{sugestoes\_categoria\_marketplace}), recuperando as categorias confirmadas anteriormente em compras de \textit{marketplace}, e armazena o resultado no estado do agente; o fluxo segue para o nó \texttt{buscar\_por\_similaridade}, executado para toda transação. Na versão avaliada, a consulta é filtrada apenas pelo usuário (não pelo estabelecimento) e o histórico recuperado \textbf{não é injetado} no \textit{prompt} do nó \texttt{classificar\_com\_llm}: o enriquecimento do contexto por histórico está previsto no estado tipado e no esquema do banco, mas não chega a influenciar a decisão -- lacuna registrada na Seção~\ref{sec:concl_limitacoes} e relevante para a resposta à Q4;
	\item \textbf{Nó \texttt{buscar\_por\_similaridade}:} Executa a busca por similaridade de cosseno na coleção vetorial do pgvector ($k = 3$ vizinhos), utilizando os \textit{embeddings} gerados pelo modelo \modelemb{} sobre a descrição normalizada. Se a melhor correspondência apresentar escore de similaridade $\ge$~0,85, a transação é considerada classificada com alta certeza pelo ramo vetorial e roteada diretamente para o nó \texttt{salvar\_resultado}, sem passar pelo ramo LLM nem pelo nó \texttt{salvar\_vectorstore} (a descrição já tem correspondente próximo na base vetorial). Os vizinhos recuperados são usados exclusivamente para essa decisão de roteamento: eles \textbf{não são repassados ao modelo de linguagem} como exemplos, de modo que o ramo LLM opera em regime \textit{zero-shot} (Seção~\ref{sec:fund_embeddings_rag});
	\item \textbf{Nó \texttt{classificar\_com\_llm}:} Caso a busca vetorial resulte em confiança inferior a 0,85, aciona-se o modelo de linguagem (LLM). O nó injeta um \textit{System Prompt} dinâmico contendo: (i) metadados da origem do arquivo; (ii) regras de decodificação de terminais de captura e adquirentes; (iii) diretrizes antialucinação; (iv) cadeia de pensamento obrigatória gerando o campo \texttt{raciocinio} antes da categoria; e (v) acesso à ferramenta \texttt{duckduckgo\_search} para pesquisar nomes comerciais obscuros -- o texto integral do \textit{prompt} encontra-se reproduzido no Apêndice~\ref{apendice:system_prompt}. A mensagem de usuário enviada ao modelo contém a \textbf{descrição bruta} da transação (``Classifique esta transação: \ldots''), não a descrição normalizada -- a normalização léxica alimenta o ramo vetorial, e o ramo LLM recebe o texto original acompanhado das regras de decodificação de \textit{gateways} do \textit{prompt}. Se o servidor de inferência rejeitar a vinculação de ferramentas (\textit{tool binding}), o nó repete a chamada sem ferramentas, de forma transparente;
	\item \textbf{Roteador unificado e nó \texttt{tools}:} Se o LLM emitir uma chamada de ferramenta (\textit{tool call}), o fluxo executa a pesquisa autônoma na web pelo \texttt{ToolNode} e retorna as evidências para o classificador;
	\item \textbf{Nó \texttt{estruturar\_saida\_llm}:} Converte a resposta em texto livre do LLM no contrato Pydantic \texttt{ClassificacaoLLM} (categoria, subcategoria, escore de confiança de 0 a 1, justificativa, lista de categorias alternativas e cadeia de raciocínio; esquema completo no Apêndice~\ref{apendice:system_prompt}) por uma cascata de três estratégias: (a) extração por expressão regular do bloco JSON que o próprio \textit{prompt} exige ao final da resposta; (b) na ausência desse bloco, nova chamada ao modelo com \texttt{with\_structured\_output} e validação Pydantic; (c) em caso de falha, uma segunda chamada de extração em texto livre com análise manual do JSON. Em todas as vias, o rótulo devolvido é \textbf{ancorado} à lista de categorias do banco por casamento exato, aproximado (\texttt{difflib}, corte 0,6) ou por substring; rótulos sem correspondência recebem a categoria residual Despesas Diversas / Outros. A implicação desse ancoramento para a leitura da matriz de confusão é discutida na Seção~\ref{sec:val_resultados_classificacao};
	\item \textbf{Nó \texttt{aguardar\_confirmacao} (HITL):} Se o escore de confiança for $\le$~0,80 ou se o LLM sinalizar \texttt{requer\_confirmacao}, o grafo executa o comando \texttt{interrupt}, pausando a linha de execução do grafo e aguardando a validação ou correção do usuário; a execução é retomada via \texttt{Command(resume=resposta)}. O ponto de verificação (\textit{checkpointer}) empregado é o \texttt{MemorySaver}, mantido em memória do processo: a retomada só é possível pelo mesmo processo que interrompeu o grafo. Em consequência, nas duas modalidades de execução desta pesquisa a interrupção é retomada programaticamente -- no teste cego aceitando a sugestão do LLM, no modo de gravação mantendo a sugestão com situação \texttt{sugerida} para revisão posterior no banco -- e nenhuma delas aguarda um revisor humano em tempo real (Seção~\ref{sec:val_mlops}). A espera efetiva por confirmação exige um \textit{checkpointer} persistente e uma interface de revisão, ambos previstos como trabalho futuro;
	\item \textbf{Nó \texttt{salvar\_vectorstore}:} Alcançado somente pelo ramo LLM: transações classificadas pelo modelo de linguagem com confiança autoavaliada $>$~0,85, ou confirmadas por humano (confiança 1,0), têm sua descrição normalizada e metadados persistidos na coleção vetorial, retroalimentando a base de conhecimento; as demais atravessam o nó sem persistência. No teste cego, esse nó é substituído por uma operação nula para impedir que o próprio teste contamine a base vetorial;
	\item \textbf{Nó \texttt{salvar\_resultado}:} Persiste a transação definitiva no PostgreSQL com integridade referencial, registrando a categoria, o método de classificação utilizado (\texttt{vetorial}, \texttt{llm}, \texttt{humano}), o escore de confiança, a \textit{flag} de confirmação pendente e a situação de revisão. A tabela de registro por inferência (\texttt{log\_classificacoes}, com \textit{prompt}, resposta bruta e tempo de processamento) está modelada no esquema, mas \textbf{não é alimentada} por este nó na versão avaliada; a rastreabilidade por transação limita-se, portanto, às colunas de método e confiança de \texttt{transacoes}, e a rastreabilidade por experimento ao MLflow (Seção~\ref{sec:val_mlops}).
\end{enumerate}

Para eliminar qualquer ambiguidade da descrição em prosa, o Quadro~\ref{quadro:arestas_agente} enumera as arestas do grafo exatamente como codificadas (\codepath{notebooks/experimentos/05_agente.py}, função \texttt{build\_graph}): cada linha é uma aresta (origem, condição, destino); arestas sem condição são incondicionais.

\begin{quadro}[htb]
	\centering
	\caption{Arestas do grafo de estados do agente (origem, condição de roteamento, destino).}
	\label{quadro:arestas_agente}
	\resizebox{\textwidth}{!}{%
	\begin{tabular}{lll}
		\toprule
		\textbf{Origem} & \textbf{Condição} & \textbf{Destino} \\
		\midrule
		\texttt{normalizar} & \texttt{eh\_marketplace} verdadeiro & \texttt{buscar\_historico\_marketplace} \\
		\texttt{normalizar} & caso contrário & \texttt{buscar\_por\_similaridade} \\
		\texttt{buscar\_historico\_marketplace} & -- & \texttt{buscar\_por\_similaridade} \\
		\texttt{buscar\_por\_similaridade} & similaridade $\ge$~0,85 & \texttt{salvar\_resultado} \\
		\texttt{buscar\_por\_similaridade} & caso contrário & \texttt{classificar\_com\_llm} \\
		\texttt{classificar\_com\_llm} & o LLM emite chamada de ferramenta & \texttt{tools} \\
		\texttt{classificar\_com\_llm} & caso contrário & \texttt{estruturar\_saida\_llm} \\
		\texttt{tools} & -- & \texttt{classificar\_com\_llm} \\
		\texttt{estruturar\_saida\_llm} & \texttt{requer\_confirmacao} ou confiança $\le$~0,80 & \texttt{aguardar\_confirmacao} \\
		\texttt{estruturar\_saida\_llm} & caso contrário & \texttt{salvar\_vectorstore} \\
		\texttt{aguardar\_confirmacao} & -- & \texttt{salvar\_vectorstore} \\
		\texttt{salvar\_vectorstore} & -- (persiste somente se confiança $>$~0,85) & \texttt{salvar\_resultado} \\
		\texttt{salvar\_resultado} & -- & fim \\
		\bottomrule
	\end{tabular}}
	\fonte{Elaborado pelo autor (\the\year), a partir de \codepath{notebooks/experimentos/05_agente.py}.}
\end{quadro}

A ferramenta de busca autônoma na web (\texttt{duckduckgo\_search}), descrita nos itens acima como parte da engenharia do agente, é habilitada por variável de ambiente (\texttt{AGENTE\_WEB}). A configuração do agente efetivamente submetida ao teste cego -- e as diferenças entre ela e a arquitetura aqui descrita -- é consolidada na Seção~\ref{sec:val_config_agente}.

Os três limiares numéricos que governam o roteamento do agente desempenham papéis distintos e não concorrentes: similaridade de cosseno $\ge$~0,85 decide a classificação direta pelo ramo vetorial (nó \texttt{buscar\_por\_similaridade}); confiança $\le$~0,80 é o único critério operacional que efetivamente aciona a intervenção humana (nó \texttt{aguardar\_confirmacao}); e confiança $\le$~0,70 é um piso interno que o LLM aplica à sua própria autoavaliação ao sinalizar estabelecimentos totalmente opacos no \textit{prompt} do nó \texttt{classificar\_com\_llm}, não um critério de roteamento independente. Como 0,70 $<$ 0,80, toda transação autoavaliada pelo LLM nesse piso já está necessariamente contida no conjunto encaminhado à curadoria humana, de modo que os dois valores não competem pela mesma decisão. Os três valores são \textbf{escolhas heurísticas do autor}: valores redondos que demarcam, respectivamente, alta confiança, incerteza relevante e baixa confiança, sem respaldo em uma fonte específica. Não foram submetidos a calibração formal nem a análise de sensibilidade sobre o conjunto de validação de 176 transações (Seção~\ref{sec:val_dataset_classificacao}), ainda que esse conjunto tenha sido reservado com essa finalidade. A ausência dessa calibração é registrada como limitação do trabalho (Seção~\ref{sec:concl_limitacoes}).

\section{Conformidade com a LGPD e segurança da informação financeira}
\label{sec:sol_lgpd}

O processamento de dados transacionais bancários envolve informações altamente sensíveis que revelam hábitos de vida, localização geográfica, saúde e patrimônio dos cidadãos. A arquitetura proposta foi desenhada em observância à Lei Geral de Proteção de Dados Pessoais (LGPD -- Lei nº 13.709/2018) \cite{brasil2018lgpd}, estruturando-se sobre quatro salvaguardas:

\begin{enumerate}
	\item \textbf{Princípio da necessidade (Art. 6º, III):} A esteira limita o tratamento ao mínimo necessário à finalidade de categorização orçamentária. Números completos de cartões de crédito, códigos de segurança (CVV), endereços residenciais e senhas não são extraídos dos documentos, preservando-se apenas a descrição textual, a data, o valor da compra e os quatro dígitos finais do cartão, necessários para distinguir cartões adicionais e virtuais de uma mesma fatura;
	\item \textbf{Pseudonimização do titular (Art. 13, § 4º):} Nenhum CPF, documento de identidade ou dado de cartão do titular é armazenado. O cadastro de usuários restringe-se ao mínimo necessário à autenticação (identificador numérico, nome, e-mail e situação), e todas as tabelas transacionais referenciam o titular exclusivamente por esse identificador numérico interno -- de modo que faturas, transações e histórico de classificação são \textbf{pseudonimizados}, no sentido legal do termo: só podem ser associados a um indivíduo com o uso de informação mantida separadamente (o cadastro e o servidor de identidade Logto, item 4). Não se trata, portanto, de anonimização, e a lei continua aplicável a esses dados;
	\item \textbf{Soberania de processamento em infraestrutura local (privacidade por design):} Ao adotar extração documental local, modelos locais de \textit{embeddings} e de linguagem e contêineres executados na infraestrutura do próprio usuário, a solução reduz substancialmente o risco de vazamento de extratos bancários para servidores de terceiros ou provedores comerciais proprietários de IA, eliminando essa via específica de exposição -- ainda que não elimine outras superfícies de risco (por exemplo, o próprio ambiente local do usuário). A única exceção é a ferramenta opcional de busca na web, que envia à internet o nome do estabelecimento (nunca valores, datas ou dados do titular);
	\item \textbf{Autenticação e rastreabilidade (Art. 46):} A aplicação \textit{web} autentica o usuário pelo servidor Logto via protocolo OpenID Connect (OIDC). Na versão avaliada, contudo, as APIs de ingestão e de classificação não validam o \textit{token} de acesso, e o controle por papéis (RBAC) não está implementado -- a segurança de acesso aos dados depende, por ora, do isolamento de rede da infraestrutura local. A rastreabilidade das decisões é atendida em dois níveis: por transação, pelas colunas de método, confiança e situação de revisão da tabela \texttt{transacoes}; e por experimento, pelo MLflow. A tabela \texttt{log\_classificacoes}, modelada para registrar cada inferência com \textit{prompt} e resposta, não é alimentada nesta versão (Seção~\ref{sec:sol_agente_inteligente}).
\end{enumerate}

\textbf{Base legal do tratamento.} O tratamento dos dados dos dois portadores de cartão realizou-se para fins exclusivamente acadêmicos, hipótese em que a LGPD dispensa parte de suas exigências (Art. 4º, II, ``b''), mas preserva a aplicação de seus princípios e a necessidade de uma base legal (Arts. 7º e 11). A base legal adotada é o \textbf{consentimento} (Art. 7º, I): o autor é um dos portadores, e o segundo portador consentiu expressamente com o uso de suas faturas nesta pesquisa. Descrições de estabelecimento, datas e valores não constituem dados pessoais sensíveis (Art. 5º, II); observa-se apenas que a categoria Saúde, inferida a partir do nome do estabelecimento, pode ser lida como indício indireto sobre a saúde do titular, razão adicional para que o acervo não seja distribuído (Apêndice~\ref{apendice:repositorio}).

\textbf{Verificação empírica.} Para que os princípios acima não fiquem apenas como asserção, a base de produção do agente foi inspecionada por consulta somente leitura (\codepath{notebooks/experimentos/10_verificacao_lgpd.py}), que conta, nas colunas textuais persistidas (3.025 transações e 488 estabelecimentos, nos campos de descrição original e normalizada, além das seis colunas textuais do registro de classificações), as ocorrências de quatro padrões de dado sensível: número de cartão (13 a 19 dígitos, com ou sem separadores), CPF (com ou sem máscara), código de verificação rotulado (``cvv''/``cvc'' seguido de dígitos) e endereço de e-mail. O resultado foi de \textbf{zero ocorrências} para CPF, CVV e e-mail em todas as colunas. O padrão de número de cartão, deliberadamente amplo, apontou quatro cadeias candidatas, todas a mesma descrição de estabelecimento (um lançamento parcelado do \textit{marketplace} Shopee e seu registro na base de conhecimento) contendo um identificador numérico de 15 dígitos que não satisfaz o algoritmo de Luhn \cite{luhn1960patent}, dígito verificador dos números de cartão -- isto é, um número de pedido do \textit{marketplace}, e não um PAN de cartão (\codepath{notebooks/experimentos/10b_inspecao_lgpd.py}, que imprime os candidatos com os dígitos mascarados). A mesma consulta confirmou que a tabela de usuários não possui coluna de CPF e que nenhum de seus campos textuais contém padrão de CPF, e que as seis colunas textuais de \texttt{log\_classificacoes} estão vazias -- consistente com o fato de essa tabela não ser alimentada nesta versão. A verificação converte a asserção de necessidade em evidência sobre o estado efetivo da base; ela não substitui, contudo, uma auditoria da camada bruta (\textit{bucket} Bronze), em que os PDFs originais das faturas -- que naturalmente contêm nome do titular e final do cartão -- são preservados por desígnio. Registra-se, por fim, que as credenciais de serviço da infraestrutura local (banco, armazenamento de objetos e servidor de identidade) constam em claro no repositório e nos arquivos de configuração; sua externalização para variáveis de ambiente é uma correção de segurança pendente, sem impacto nos resultados.

\section{Estratégia de garantia de qualidade e testes}
\label{sec:sol_testes}

Para um sistema distribuído e poliglota de natureza exploratória, a garantia de qualidade apoiou-se em três frentes, de maturidade desigual:

\begin{itemize}
	\item \textbf{Testes de regressão de modelo:} A esteira experimental reprodutível (\codepath{notebooks/experimentos/}, etapas 01 a 13) fixa a semente aleatória e um conjunto de teste cego particionado por descrição normalizada (Seção~\ref{sec:val_dataset_classificacao}), permitindo detectar regressões nas métricas preditivas a cada alteração de código, \textit{prompt} ou hiperparâmetro, com os resultados versionados no MLflow. Esta é a frente efetivamente automatizada e exercitada;
	\item \textbf{Validação de integração assistida:} A comunicação de ponta a ponta entre a API de ingestão, o barramento Kafka, o armazenamento de objetos MinIO e o serviço consumidor de extração foi exercitada de forma assistida em ambiente Docker Compose, confirmando a publicação e o consumo de eventos e a persistência dos arquivos brutos e das transações extraídas;
	\item \textbf{Testes unitários:} O ambiente Python está preparado para \textit{pytest} (com suporte a execução assíncrona), mas a suíte de testes unitários das funções determinísticas de pré-processamento --- saneamento léxico com \textit{Regex}, analisadores (\textit{parsers}) de datas e valores monetários, resolução de \textit{aliases} comerciais e contratos Pydantic --- \textbf{não foi escrita} nesta versão; o serviço de ingestão em Java conta apenas com a verificação de carga do contexto \textit{Spring Boot} via \textit{JUnit~5}. A ausência dessa suíte explica, em parte, que defeitos determinísticos do agente (como a chave de \textit{marketplace} não casada, Apêndice~\ref{apendice:normalizacao}) não tenham sido detectados antes da execução dos experimentos.
\end{itemize}

A suíte de testes unitários e a automação plena da esteira de testes (integração com \textit{Testcontainers} e testes de carga do barramento) são registradas como trabalho futuro no Capítulo~\ref{Cap:Conclusoes}.

\section{Considerações finais do capítulo}
\label{sec:sol_consideracoes_finais}

Este capítulo formalizou as questões de pesquisa e a hipótese central, com seus pares de hipóteses nula e alternativa, e descreveu a solução proposta em suas decisões de projeto mais consequentes: a extração nativa da camada de texto do PDF com preservação do leiaute tabular, em vez de rasterização seguida de OCR; o roteamento por limiar de similaridade, que reserva o modelo de linguagem aos estabelecimentos não reconhecidos pela base vetorial; a execução de todos os componentes -- extração, \textit{embeddings} e LLM -- em infraestrutura local, sob as salvaguardas da LGPD; e o particionamento do acervo por descrição normalizada, que impede a repetição literal de lançamentos entre treino e teste. Registrou também, com a mesma precisão, as lacunas de integração da versão avaliada: a extração óptica não acoplada, a ferramenta de busca e o histórico de \textit{marketplace} sem efeito sobre a decisão, a confirmação humana simulada, o registro por inferência não alimentado e a ausência de suíte de testes unitários. O capítulo seguinte submete essa versão ao protocolo experimental: caracteriza os conjuntos de dados, aplica o particionamento e discute os seus limites, compara as cinco arquiteturas pelas métricas e pelos testes estatísticos formalizados no Capítulo~\ref{Cap:Fund_Teo} e verifica as duas sub-hipóteses.






